\section{Composite mass properties and the two cache gates}\label{sec:caching}

Every physical question the simulator asks of the rocket --- what does it weigh, where is its
center of mass, what is its inertia tensor --- is a \emph{composite} question: a mass-weighted
aggregation over the part tree of \cref{sec:tree}, evaluated at a simulation time $t$ because the
motor's mass falls while it burns. These aggregations are answered by \code{PartNode}
(\srcf{core/model/PartsModel.h}), and any node answers them for its own sub-tree, so a booster
stack and a full rocket use the same code path. The design problem is that the same queries sit on
two wildly different clocks. The integrator asks thousands of times per flight, at times chosen by
an adaptive step controller that revisits and discards trial times freely. The designer asks after
each structural edit, at human cadence. A single ``recompute when dirty'' flag cannot serve both,
because during a flight nothing ever marks the tree dirty --- time itself is the thing that
changes.

The resolution is two independent caches per node, each invalidated by exactly the events that can
change its output. The \emph{placement cache} (resolved poses plus the overlap verdict of
\cref{sec:placement}) is keyed on structure: it rebuilds only when a part is attached, detached, or
re-linked, and never during a flight, because geometry is invariant under mass loss. The
\emph{composite cache} (CM and inertia tensor) is keyed on the one physical quantity that does
change in flight: the composite mass itself. This section walks through the quantities, the two
gates, the two-pass mathematics behind the cached values, the aero and reference-area aggregates
that share the same resolve, and what the flight loop actually pays.

\subsection{One live sum, two cached tensors}\label{sec:caching-quantities}

\code{PartNode} exposes three time-aware composite quantities
(\src{core/model/PartsModel.h}{70}):

\begin{itemize}
\item \code{compositeMass(t)} --- a deliberately \emph{uncached} recursive sum of
      \code{Part::getMass(t)} over the sub-tree (\src{core/model/PartsModel.cpp}{102}). It is both
      the ODE divisor ($a = F/m$) and the key of the mass-delta gate.
\item \code{compositeCm(t)} --- the mass-weighted center of mass in the sub-tree root's
      fore-plane frame, cached per node behind the gate.
\item \code{compositeI(t)} --- the full mass-weighted inertia tensor
      ($\unit{kg\,m^2}$) about the composite CM, cached behind the same gate.
\end{itemize}

The asymmetry is the point. The CM and tensor are worth caching: they require the resolved
placements, an $O(N)$ weighted average, and $N$ parallel-axis shifts of $3\times3$ tensors. The
mass sum is $N$ double additions with no tensor work at all --- cheaper than the Eigen expression
it feeds --- so caching it would buy nothing. But the deeper reason it must stay live is
structural:

\begin{keyinsight}[The live sum is the tree's clock]
Nothing notifies the part tree that simulation time advanced. No verb fires, no flag flips; the
only place $t$ enters the mass properties is \code{Part::getMass(t)}. The gate therefore needs
one authoritative, freshly-evaluated ``mass now'' to compare against its stamp. If
\code{compositeMass} were itself cached, the gate would be handed its own stamp back:
\code{mNow} would equal \cppinline|builtAtMass_| by construction, the comparison could never fail,
and the CM and tensor would silently stop tracking the burn. The live sum is not an optimization
left on the table --- it is the input that makes the cache correct.
\end{keyinsight}

\subsection{The mass-delta gate}\label{sec:caching-massgate}

Both cached accessors funnel through \code{ensureCompositeCache} (\cref{lst:caching-ensure}). The
gate condition is two-legged: \cppinline|massDirty_| is the structural leg, raised by every
\code{PartsModel} verb whose edit can move mass (\cref{sec:tree}); the mass comparison is the
flight leg, and it is an \emph{exact} floating-point equality.

\begin{listing}[htbp]
\begin{minted}{cpp}
double PartNode::compositeMass(double t) const
{
    // ...
    double m = part_->getMass(t);
    for(const auto& child : children_)
    {
        m += child->compositeMass(t);
    }
    return m;
}

void PartNode::ensureCompositeCache(double t) const
{
    // ...
    const double mNow = compositeMass(t);
    if(massDirty_ || mNow != builtAtMass_)
    {
        const Composite c = computeCompositeAt(t);
        compositeCm_ = c.cm;
        compositeI_  = c.inertia;
        builtAtMass_ = mNow;
        massDirty_   = false;
    }
}
\end{minted}
\caption{The live sum and the mass-delta gate (\srccap{core/model/PartsModel.cpp}{102}). Stamps
are written only after \code{computeCompositeAt} succeeds, so a failed placement solve is never
absorbed by the cache (\cref{sec:caching-composite}).}\label{lst:caching-ensure}
\end{listing}

\cppinline|builtAtMass_| is initialized to \cppinline|quiet_NaN()|
(\src{core/model/PartsModel.h}{116}). NaN compares unequal to everything, including itself, so the
very first read builds unconditionally --- the sentinel makes the gate self-priming without a
separate ``ever built'' flag.

Exact equality looks fragile next to floating point folklore, but here it is precisely correct,
because both sides of the comparison are produced by the \emph{same} deterministic computation:

\begin{itemize}
\item \textbf{Post-burnout the cache freezes.} After the propellant is spent,
      \code{MotorModel::getMass} returns the stored casing mass verbatim for every $t$
      (\src{core/model/MotorModel.cpp}{65}), and every other part's mass ignores $t$ entirely. The
      recursive sum runs over a fixed tree in fixed child order, adding identical doubles in
      identical sequence, so it reproduces \emph{bit-for-bit} on every read. The gate compares
      equal, and the coast and descent phases --- the majority of most flights --- never rebuild.
\item \textbf{During a burn the cache tracks.} Between recorded steps the interpolated mass curve
      moves, the sum differs from the stamp, and the gate rebuilds --- once per recorded step,
      as it must, because the CM and tensor genuinely moved.
\item \textbf{The key is the physics, not the schedule.} In flight the motor is the only
      time-varying contributor, so the total mass determines the composite state: equal totals at
      two different times imply identical CM and tensor, and a rebuild would produce the value
      already stored. Concurrent mass \emph{edits} that could confound the total are impossible
      mid-read --- they route through verbs that raise \cppinline|massDirty_|, the gate's other
      leg.
\end{itemize}

\begin{rejected}[Rejected --- tolerance comparison, and time as the key]
An epsilon comparison (\,$|m_{\text{now}} - m_{\text{built}}| < \varepsilon$\,) would freeze the
cache during any burn slow enough that per-step mass loss drops below $\varepsilon$ --- a silently
stale CM on exactly the flights where sustainer motors matter. Exact \code{==} has no false-freeze
mode: the only way the sum repeats is that the mass genuinely did not change.

Keying the cache on time (``rebuilt at $t = 3.117$'') fails against the adaptive integrator.
RK45 evaluates six Fehlberg stages per attempted step at intermediate node times
(\src{core/sim/RK45Solver.h}{86}) and discards the whole attempt when the error estimate is out of
tolerance; trial times are visited, abandoned, and never exactly revisited. A monotonic time stamp
would be poisoned by a rejected attempt that probed \emph{ahead} of the eventually accepted step.
The mass key is indifferent to all of this: a rebuild at a rejected trial mass is simply
overwritten by the next read, and after burnout no time query of any kind can move the key.
\end{rejected}

\Cref{fig:caching-gate} traces the full read path, including the failure branch discussed in
\cref{sec:caching-composite}.

\tikzfig{activity-composite-gate}{The composite read path. The live mass sum is compared
exact-\code{==} against the stamp (NaN sentinel on first read); a rebuild runs the placement gate
first, refuses a failed overlap verdict by throwing, and stamps the cache only after
success.}{fig:caching-gate}

\subsection{The placement gate}\label{sec:caching-placementgate}

The second cache holds the resolved placements --- one \code{Placed} pose per part of the
sub-tree, in deterministic depth-first attachment order --- plus the overlap-sweep verdict, both
produced by the resolver machinery of \cref{sec:placement}. \code{ensurePlacementCache}
(\src{core/model/PartsModel.cpp}{144}) rebuilds both iff \cppinline|placementDirty_|, which is
raised only by structural and link edits: attach, detach, \code{installRoot}, \code{setLink}, and
\code{mutatePart} under the \code{Geometry} hint (\src{core/model/PartsModel.h}{139}). A mass or
inertia edit --- \code{setPartMass}, \code{setPartInertia}, a motor swap, or the burn itself ---
never touches it, because none of those can move a seam. The expensive $O(N \log N)$ envelope
sweep therefore runs exactly once per edit, at human cadence, and the flight loop reuses its
output thousands of times.

\begin{designrule}[Two gates, two cadences]
Each cache is invalidated by exactly the events that can change its output, and by nothing else.
The placement cache ticks at \emph{editing} cadence: resolve and sweep once per structural or link
edit, never during a flight, because geometry is invariant under mass loss. The composite cache
ticks at \emph{simulation} cadence: rebuild once per recorded step while mass is changing, then
never again. Burning propellant cannot trigger a geometry solve; dragging a fin cannot be missed
by the mass gate.
\end{designrule}

The last clause needs one supporting rule. A link edit can be mass-neutral --- sliding a nose cone
forward changes no part's mass, so the mass gate, keyed on the total, is structurally blind to it.
\code{setLink} therefore dirties \emph{both} chains (\src{core/model/PartsModel.cpp}{443}): the
placement chain because poses moved, and the mass chain because the composite CM and tensor depend
on those poses. The two gates overlap by design wherever one alone would be blind. Dirty marking
itself is an $O(\text{depth})$ walk up the parent pointers (\src{core/model/PartsModel.h}{103}),
so an edit anywhere invalidates every enclosing sub-assembly, and sibling sub-trees keep their
caches untouched.

\subsection[computeCompositeAt: two passes over one resolve]{\code{computeCompositeAt}: two passes over one resolve}\label{sec:caching-composite}

A gate-triggered rebuild calls \code{computeCompositeAt(t)}
(\src{core/model/PartsModel.cpp}{171}), which re-weights the \emph{cached} geometry by the
\emph{current} masses --- the placement gate has already guaranteed the resolved poses are
current. The math is two passes over \code{resolvedPlacements}.

\textbf{Pass 1 --- mass-weighted CM.} Each part $i$ stores its CM relative to its component
middle (\cref{sec:leaf}), so its CM in local coordinates, its position in the sub-tree-root frame,
and the composite CM are
\begin{equation}
\mathbf{c}^{\mathrm{loc}}_i =
\begin{pmatrix} c_{i,x}\\[2pt] c_{i,y}\\[2pt] -\tfrac{L_i}{2} + c_{i,z} \end{pmatrix},
\qquad
\mathbf{r}_i = \mathbf{o}_i + q_i\,\mathbf{c}^{\mathrm{loc}}_i,
\qquad
\bar{\mathbf{r}}(t) = \frac{\sum_i m_i(t)\,\mathbf{r}_i}{\sum_i m_i(t)},
\label{eq:caching-cm}
\end{equation}
where $\mathbf{o}_i$ and $q_i$ are the part's resolved fore-plane origin and orientation
(identity in 3-DOF), and $L_i$ its length --- a part occupies $z \in [-L_i, 0]$ in its own frame,
\zfwd.

\textbf{Pass 2 --- parallel-axis tensor.} Each part's per-unit-mass geometric tensor
$\mathbf{I}^{\mathrm{geo}}_i$ ($\unit{m^2}$, about its own CM; \cref{sec:leaf}) is weighted to a
real tensor and shifted to the composite CM by the parallel-axis theorem:
\begin{equation}
\mathbf{I}(t) \;=\; \sum_i m_i(t)\,\mathbf{I}^{\mathrm{geo}}_i
\;+\; m_i(t)\left[(\mathbf{d}_i\!\cdot\!\mathbf{d}_i)\,\mathbf{E}
- \mathbf{d}_i\,\mathbf{d}_i^{\mathsf{T}}\right],
\qquad
\mathbf{d}_i = \mathbf{r}_i - \bar{\mathbf{r}}(t),
\label{eq:caching-pat}
\end{equation}
with $\mathbf{E}$ the identity. The displacement term is even in $\mathbf{d}_i$, so the sign
convention of the offset never matters. \Cref{lst:caching-passes} is the code, line for line the
same shape as the equations.

\begin{listing}[htbp]
\begin{minted}{cpp}
Matrix3 parallelAxisTerm(const Vector3& d)
{
    return d.dot(d) * Matrix3::Identity() - d * d.transpose();
}
// ...
    // Pass 1: mass-weighted composite CM.
    double  m        = 0.0;
    Vector3 weighted = Vector3::Zero();
    for(const part::Placed& pl : resolved_)
    {
        const double  pm       = pl.part->getMass(t);
        const Vector3 cmOffset = pl.part->getCenterMassOffset();
        const Vector3 cmLocal(cmOffset.x(), cmOffset.y(),
                                     -pl.part->getLength() / 2.0 + cmOffset.z());
        const Vector3 cmInRoot = pl.pose.origin + pl.pose.orient * cmLocal;
        m        += pm;
        weighted += pm * cmInRoot;
        parts.push_back(Contribution{cmInRoot, pl.part, pm});
    }
    // Guard the divide: a fully massless sub-tree has no meaningful CM, so leave it at the origin.
    Vector3 cm = Vector3::Zero();
    if(m > 0.0)
    {
        cm = weighted / m;
    }

    // Pass 2: inertia about the composite CM. Shift each part's own mass-weighted tensor to cm via
    // the parallel-axis theorem. 6-DOF adds the R*I*R^T rotation (identity today) here.
    Matrix3 I = Matrix3::Zero();
    for(const Contribution& c : parts)
    {
        const Vector3 d = c.cmInRoot - cm;
        I += c.mass * c.part->getI() + c.mass * parallelAxisTerm(d);
    }
\end{minted}
\caption{The two-pass composite walk (\srccap{core/model/PartsModel.cpp}{171}, elided). The
massless-subtree guard keeps a zero-mass assembly --- legal while sketching a design --- from
producing a $0/0$ CM; it reads as the origin instead.}\label{lst:caching-passes}
\end{listing}

Before either pass runs, the walk consults the cached overlap verdict, and this is where the
design refuses to be quietly wrong:

\begin{hardfail}[A self-intersecting design cannot be flown]
When the envelope sweep of \cref{sec:placement} finds interpenetrating parts,
\code{computeCompositeAt} throws \cppinline|std::runtime_error| carrying the first located
diagnostic (\src{core/model/PartsModel.cpp}{184}) rather than integrating overlapping volumes into
a plausible-looking but physically meaningless tensor. Because the stamps in
\cref{lst:caching-ensure} land only after a successful build, \cppinline|massDirty_| stays raised
and \emph{every} subsequent \code{compositeCm}/\code{compositeI} read re-throws --- the failure
cannot be absorbed by asking twice. The repeat throws are cheap: \cppinline|placementDirty_| was
cleared when the verdict was cached, so no re-sweep runs; the verdict itself is the cached
diagnosis. The condition clears only the way it arose --- through a structural edit that re-opens
the placement gate and produces a clean verdict.
\end{hardfail}

\subsection{Aero and area: aggregates on the same resolve}\label{sec:caching-aero}

Two further aggregates fold over the same cached resolve. \code{compositeAero(refArea)}
(\src{core/model/PartsModel.cpp}{224}) assembles the whole-rocket Barrowman profile. Each part
reports an \code{AeroComponent} whose CP station $x_{cp,i}$ is measured from \emph{that part's own
CM} (\srcf{core/model/Aero.h}), stored as the weighted moment $C_{N\alpha,i}\,x_{cp,i}$ so that a
zero-lift body drops out of the average with no special case. The composite walk re-expresses
every part onto the shared sub-tree-root datum by adding back the part's CM station:
\begin{equation}
z_{cm,i} = o_{i,z} - \tfrac{L_i}{2} + c_{i,z},
\qquad
x_{cp} = \frac{\sum_i C_{N\alpha,i}\left(x_{cp,i} + z_{cm,i}\right)}
              {\sum_i C_{N\alpha,i}}.
\label{eq:caching-datum}
\end{equation}
The added $z_{cm,i}$ exactly cancels the per-part CM datum that $x_{cp,i}$ was measured from, so
the composite CP depends on external shape alone --- and it now lives on the same fore-plane datum
as $\bar{\mathbf{r}}$ from \cref{eq:caching-cm}. That shared datum is the payoff: the static
margin is literally \code{profile.cp() - cg}, one subtraction with no frame conversion. The
profile is re-derived on each call rather than cached; it sits on no flight-loop path, and its
only nontrivial input --- the resolved placements --- is already cached behind the placement gate.

\code{maxFrontalReferenceArea} (\src{core/model/PartsModel.cpp}{243}) supplies the drag reference
area as a \emph{maximum}, not a sum: the largest single frontal disc over the sub-tree, the
Barrowman convention. Summing would double-count --- a nose cone in front of a body tube of the
same caliber presents one disc to the flow, not two --- and a fin set deliberately reports only
the disc of the body it mounts on (\src{core/model/parts/FinSet.h}{71}), so adding fins never
inflates the reference area. \code{RocketModel} exposes it as
\code{deriveReferenceAreaFromGeometry()} for front ends (\src{core/model/RocketModel.cpp}{44});
the drag term itself always reads the stored \code{referenceArea} scalar, which
\code{setReferenceArea} marks as a manual override (\src{core/model/RocketModel.h}{69}).

\subsection{What the flight loop pays}\label{sec:caching-costs}

The gates exist so that the composite machinery costs the integrator almost nothing.
\Cref{tab:caching-costs} itemizes the work by event; \cref{fig:caching-timeline} shows the
resulting cadence over a flight.

\begin{table}[htbp]
\centering
\small
% ragged-right wrapping columns: the long event labels and mono src citations don't justify
\begin{tabularx}{\linewidth}{@{}>{\raggedright\arraybackslash}p{0.26\linewidth}>{\raggedright\arraybackslash}X>{\raggedright\arraybackslash}X@{}}
\toprule
\textbf{Event} & \textbf{Work} & \textbf{Cache traffic} \\
\midrule
Integrator stage (RK4: 4 per step; RKF45: 6 per attempt, re-run on rejection)
  & Two live \code{compositeMass} sums: the gravity weight inside \code{getForces}
    (\src{core/model/RocketModel.cpp}{73}) and the $F/m$ divide in the ODE
    (\src{core/sim/Propagator.cpp}{41}). No tensor math.
  & None. The stage path never reads the cached CM or tensor. \\
\addlinespace
Accepted, recorded step
  & \code{writeMassProperties} snapshots mass, CG, and tensor into the state history
    (\src{core/model/RocketModel.cpp}{49}, called at
    \src{core/sim/Propagator.cpp}{132}): one more live sum plus the two gated reads.
  & At most one gate rebuild while the motor burns; exactly zero after burnout. \\
\addlinespace
Structural / link edit
  & Dirty walk to the root, $O(\text{depth})$ flag writes.
  & Next read: one resolve + envelope sweep, then one composite rebuild. \\
\bottomrule
\end{tabularx}
\caption{Per-event cost of the composite machinery. The hot path touches only the mass-only
recursive sum; the tensor work runs once per recorded step during a burn and not at all
afterward.}\label{tab:caching-costs}
\end{table}

The RKF45 interaction deserves emphasis because it is where a naive cache design breaks. A
rejected attempt evaluates all six stages at trial times and throws the results away; as wired,
those stages touch nothing but the live sum, so there is nothing to poison. And even a future
consumer that read \code{compositeI} mid-stage would be safe, because the gate keys on the mass
\emph{value}: a rebuild triggered at a discarded trial mass is just a valid cache entry for that
mass, corrected by the next read that sees a different one. Correctness is independent of the step
controller's accept/reject sequence --- a property a time-keyed cache cannot offer.

\tikzfig{timeline-burn}{Gate cadence over one flight. During the burn the live sum differs from
the stamp at every recorded step, so the gate opens once per step and the CM and tensor track the
mass loss. From burnout on, the sum reproduces bit-for-bit, exact \code{==} holds, and every read
is a frozen-cache hit.}{fig:caching-timeline}

\subsection{The single-writer contract}\label{sec:caching-singlewriter}

All of the above state --- both dirty flags, the stamp, the cached CM and tensor, the resolved
placements and verdict --- is \code{mutable}, written behind \code{const} accessors. A
\code{const PartNode\&} is therefore \emph{not} a promise of physical immutability; it is a
promise of logical immutability under the model's single-writer discipline
(\cref{sec:philosophy}). Concurrent access to a \code{PartsModel} --- even all-\code{const},
even two readers racing the same \code{compositeI(t)} call --- is a data race, because either
read may be the one that rebuilds. The contract is stated at the class boundary
(\src{core/model/PartsModel.h}{36}) rather than paid for with per-node locks: QtRocket's
consumers (the GUI thread, the CLI loop, one propagator at a time) are serialized by
construction, and a mutex around every composite read would tax the hot path of
\cref{tab:caching-costs} to defend against a caller the architecture already forbids.
